\textbf{Phần 1. Trắc nghiệm lựa chọn}
\setcounterpageref{ex}{0}

\begin{ex} %[3D3V35]
Tính: $35\text{ mm} \times 4 + 60\text{ mm}$. Kết quả theo đơn vị xăng-ti-mét là
\choice
{$200\text{ cm}$}
{$2\text{ cm}$}
{\True $20\text{ cm}$}
{$18\text{ cm}$}
\loigiai{
$35 \times 4 = 140\text{ mm}$; $140 + 60 = 200\text{ mm} = 20\text{ cm}$.
\textbf{Đáp án của câu hỏi là C}
}
\end{ex}

\begin{ex} %[3D3V35]
Có ba bao thóc. Bao thứ nhất nặng $350\text{ g}$, bao thứ hai nặng gấp đôi bao thứ nhất, bao thứ ba nặng $450\text{ g}$. Tổng khối lượng của cả ba bao thóc là
\choice
{\True $1\text{ kg } 500\text{ g}$}
{$1\text{ kg } 400\text{ g}$}
{$1\text{ kg } 600\text{ g}$}
{$2\text{ kg}$}
\loigiai{
Bao hai nặng: $350 \times 2 = 700\text{ g}$. Tổng ba bao: $350 + 700 + 450 = 1500\text{ g} = 1\text{ kg } 500\text{ g}$.
\textbf{Đáp án của câu hỏi là A}
}
\end{ex}

\begin{ex} %[3D3V35]
Một bình chứa $1\ell$ nước. Người ta rót ra 3 cốc, mỗi cốc chứa $200\text{ ml}$ và 2 chén, mỗi chén chứa $50\text{ ml}$. Trong bình còn lại bao nhiêu mi-li-lít nước?
\choice
{$400\text{ ml}$}
{$250\text{ ml}$}
{$350\text{ ml}$}
{\True $300\text{ ml}$}
\loigiai{
Lượng nước đã rót: $200 \times 3 + 50 \times 2 = 600 + 100 = 700\text{ ml}$. Lượng nước còn lại: $1000 - 700 = 300\text{ (ml)}$.
\textbf{Đáp án của câu hỏi là D}
}
\end{ex}

\begin{ex} %[3D3V35]
Nhiệt độ nước trong cốc thứ nhất là $45^\circ\text{C}$, cốc thứ hai nóng hơn cốc thứ nhất $15^\circ\text{C}$. Nhiệt độ của cốc thứ hai là
\choice
{$55^\circ\text{C}$}
{\True $60^\circ\text{C}$}
{$65^\circ\text{C}$}
{$50^\circ\text{C}$}
\loigiai{
Nhiệt độ cốc thứ hai là: $45 + 15 = 60^\circ\text{C}$.
\textbf{Đáp án của câu hỏi là B}
}
\end{ex}

\begin{ex} %[3D3V35]
Một mảnh vải dài $2\text{ m } 40\text{ cm}$. Người ta cắt lấy $\dfrac{1}{3}$ mảnh vải để may túi. Phần vải còn lại dài bao nhiêu xăng-ti-mét?
\choice
{$80\text{ cm}$}
{$180\text{ cm}$}
{\True $160\text{ cm}$}
{$120\text{ cm}$}
\loigiai{
Đổi $2\text{ m } 40\text{ cm} = 240\text{ cm}$. Đã may túi: $240 : 3 = 80\text{ cm}$. Còn lại: $240 - 80 = 160\text{ (cm)}$.
\textbf{Đáp án của câu hỏi là C}
}
\end{ex}

\begin{ex} %[3D3V35]
Một con gà cân nặng $2\text{ kg } 200\text{ g}$, con vịt nhẹ hơn con gà $500\text{ g}$. Cả gà và vịt cân nặng là
\choice
{\True $3\text{ kg } 900\text{ g}$}
{$3\text{ kg } 700\text{ g}$}
{$4\text{ kg}$}
{$3\text{ kg } 800\text{ g}$}
\loigiai{
Con vịt nặng: $2200 - 500 = 1700\text{ g}$. Cả hai con nặng: $2200 + 1700 = 3900\text{ g} = 3\text{ kg } 900\text{ g}$.
\textbf{Đáp án của câu hỏi là A}
}
\end{ex}

\begin{ex} %[3D3V35]
Một can dầu chứa $10\ell$ dầu hỏa. Mỗi lít dầu hỏa cân nặng khoảng $800\text{ g}$. Vỏ can nặng $500\text{ g}$. Cả can dầu đó cân nặng bao nhiêu?
\choice
{$8\text{ kg}$}
{$9\text{ kg}$}
{$7\text{ kg } 500\text{ g}$}
{\True $8\text{ kg } 500\text{ g}$}
\loigiai{
Số dầu cân nặng: $800 \times 10 = 8000\text{ g} = 8\text{ kg}$. Cả can dầu nặng: $8\text{ kg} + 500\text{ g} = 8\text{ kg } 500\text{ g}$.
\textbf{Đáp án của câu hỏi là D}
}
\end{ex}

\begin{ex} %[3D3V35]
Có 4 chai mật ong, mỗi chai chứa $250\text{ ml}$. Người ta đã bán đi 3 chai. Lượng mật ong còn lại là
\choice
{$500\text{ ml}$}
{\True $250\text{ ml}$}
{$750\text{ ml}$}
{$1\ell$}
\loigiai{
Còn lại 1 chai mật ong, tức là $250\text{ ml}$.
\textbf{Đáp án của câu hỏi là B}
}
\end{ex}

\begin{ex} %[3D3V35]
Một thanh sắt dài $1\text{ m } 50\text{ mm}$ được cắt thành các đoạn dài $15\text{ mm}$. Số đoạn sắt cắt được là
\choice
{$50$ đoạn}
{$60$ đoạn}
{\True $70$ đoạn}
{$80$ đoạn}
\loigiai{
Đổi $1\text{ m } 50\text{ mm} = 1050\text{ mm}$. Số đoạn cắt được: $1050 : 15 = 70$ đoạn.
\textbf{Đáp án của câu hỏi là C}
}
\end{ex}

\begin{ex} %[3D3V35]
Người ta xếp 6 gói mì chính, mỗi gói nặng $150\text{ g}$ và 2 gói muối, mỗi gói nặng $500\text{ g}$ vào một giỏ quà. Tổng khối lượng hàng trong giỏ là
\choice
{\True $1\text{ kg } 900\text{ g}$}
{$1\text{ kg } 800\text{ g}$}
{$2\text{ kg}$}
{$1\text{ kg } 700\text{ g}$}
\loigiai{
Tổng khối lượng: $150 \times 6 + 500 \times 2 = 900 + 1000 = 1900\text{ g} = 1\text{ kg } 900\text{ g}$.
\textbf{Đáp án của câu hỏi là A}
}
\end{ex}

\begin{ex} %[3D3V35]
Tìm một số biết rằng lấy số đó nhân với $25\text{ ml}$ thì được đúng $1\ell$. Số đó là
\choice
{$25$}
{$30$}
{$50$}
{\True $40$}
\loigiai{
Đổi $1\ell = 1000\text{ ml}$. Số cần tìm là: $1000 : 25 = 40$.
\textbf{Đáp án của câu hỏi là D}
}
\end{ex}

\begin{ex} %[3D3V35]
Hôm qua nhiệt độ lúc trưa là $31^\circ\text{C}$. Hôm nay do trời mưa, nhiệt độ trưa giảm đi $5^\circ\text{C}$. Hôm nay nhiệt độ lúc trưa là
\choice
{$27^\circ\text{C}$}
{\True $26^\circ\text{C}$}
{$25^\circ\text{C}$}
{$28^\circ\text{C}$}
\loigiai{
Nhiệt độ hôm nay là: $31 - 5 = 26^\circ\text{C}$.
\textbf{Đáp án của câu hỏi là B}
}
\end{ex}

\textbf{Phần 2. Trắc nghiệm đúng sai}
\setcounterpageref{ex}{0}

\begin{ex} %[3D3V35]
Xét tính đúng sai của các nhận định tổng hợp về đại lượng:
\choiceTF
{\True Độ dài được đo bằng các đơn vị: $\text{mm}, \text{cm}, \text{dm}, \text{m}, \text{km}$.}
{Gam ($\text{g}$) là đơn vị đo dung tích chất lỏng.}
{\True $1\ell = 1000\text{ ml}$ và $1\text{ kg} = 1000\text{ g}$.}
{$100\text{ cm} = 1000\text{ mm}$ là sai.}
\loigiai{
\begin{itemize}
    \item Ý a) Đúng: Hệ đơn vị độ dài.
    \item Ý b) Sai: Sai vì gam là đơn vị đo khối lượng.
    \item Ý c) Đúng: Quan hệ tỉ lệ 1000.
    \item Ý d) Sai: Sai vì $100\text{ cm} = 1\text{ m} = 1000\text{ mm}$ là đúng.
\end{itemize}
\begin{itemize}
    \item a) \textbf{Đúng}.
    \item b) \textbf{Sai}.
    \item c) \textbf{Đúng}.
    \item d) \textbf{Sai}.
\end{itemize}
}
\end{ex}

\begin{ex} %[3D3V35]
Một hộp sữa tươi có khối lượng tịnh $1\text{ kg}$, dung tích ghi $950\text{ ml}$:
\choiceTF
{\True Khối lượng của hộp sữa được đo bằng gam hoặc ki-lô-gam.}
{\True Dung tích sữa bên trong được đo bằng mi-li-lít hoặc lít.}
{Một lít sữa tươi luôn có thể tích đúng bằng $100\text{ ml}$.}
{\True Uống hết một nửa hộp sữa tức là đã uống $475\text{ ml}$ sữa.}
\loigiai{
\begin{itemize}
    \item Ý a) Đúng: Đơn vị đo khối lượng.
    \item Ý b) Đúng: Đơn vị đo dung tích.
    \item Ý c) Sai: Sai vì $1\ell = 1000\text{ ml}$.
    \item Ý d) Đúng: $950 : 2 = 475\text{ ml}$.
\end{itemize}
\begin{itemize}
    \item a) \textbf{Đúng}.
    \item b) \textbf{Đúng}.
    \item c) \textbf{Sai}.
    \item d) \textbf{Đúng}.
\end{itemize}
}
\end{ex}

\textbf{Phần 3. Trả lời ngắn}
\setcounterpageref{ex}{0}

\begin{ex} %[3D3V35]
Tính giá trị biểu thức: $250\text{ g} \times 4 - 300\text{ g}$ theo đơn vị gam.
\shortans{700}
\loigiai{
\textbf{Đáp số:} $700$\\
$1000 - 300 = 700\text{ (g)}$.
}
\end{ex}

\begin{ex} %[3D3V35]
Đổi $3\text{ m } 5\text{ cm}$ ra đơn vị mi-li-mét.
\shortans{3050}
\loigiai{
\textbf{Đáp số:} $3050$\\
$3000 + 50 = 3050\text{ (mm)}$.
}
\end{ex}

\begin{ex} %[3D3V35]
Một bình thủy có $2\ell$ nước sôi. Người ta rót ra 5 tách trà, mỗi tách $150\text{ ml}$. Hỏi trong bình còn lại bao nhiêu mi-li-lít nước sôi?
\shortans{1250}
\loigiai{
\textbf{Đáp số:} $1250$\\
Đã rót: $150 \times 5 = 750\text{ ml}$. Còn lại: $2000 - 750 = 1250\text{ (ml)}$.
}
\end{ex}

\begin{ex} %[3D3V35]
Một cây kem que có nhiệt độ $-5^\circ\text{C}$, sau khi để ngoài trời nhiệt độ tăng thêm $8^\circ\text{C}$. Hỏi nhiệt độ cây kem lúc này là bao nhiêu độ C?
\shortans{3}
\loigiai{
\textbf{Đáp số:} $3$\\
$-5 + 8 = 3^\circ\text{C}$.
}
\end{ex}

\begin{ex} %[3D3V35]
Một quả dưa lưới nặng $1\text{ kg } 600\text{ g}$. Quả dưa hấu nặng gấp đôi quả dưa lưới. Hỏi quả dưa hấu nặng bao nhiêu gam?
\shortans{3200}
\loigiai{
\textbf{Đáp số:} $3200$\\
$1600 \times 2 = 3200\text{ (g)}$.
}
\end{ex}

\begin{ex} %[3D3V35]
Cần bao nhiêu ca nước dung tích $250\text{ ml}$ để đổ đầy một thùng chứa $5\ell$?
\shortans{20}
\loigiai{
\textbf{Đáp số:} $20$\\
$5000 : 250 = 20$ ca.
}
\end{ex}
